\documentclass{ar2rc}

\title{Neuromorphic Computing Based on Parametrically-Driven \\Oscillators and Frequency Combs}
\author{Mahadev Sunil Kumar and Adarsh Ganesan}
\journal{Applied Physics Letters}


\begin{document}

	\maketitle

	Dear Reviewers,

	Thank you for your careful and constructive reviews. We have fixed the typos as well as made other changes as described in the reviews, in the manuscript.

	\newpage

	\section{Reviewer \#1}
	\RC In Fig. 2, the title of panel (a1) should be $\Delta F = F_{avg}$, not $\Delta F = 0$.

	\AR The typo has been fixed and Fig. 1 in the manuscript has been updated.

	\section{Reviewer \#2}
	\RC Add concise physical interpretation of IPC in the main text and note that tabulated IPC values are lower-bounds limited by polynomial-order truncation.

	\AR We have added the results of the IPC tests in the main text as well as emphasized that the values are upper bounds with the IPC value of parametric resonance still increasing at the $11^{\text{th}}$ degree.

	\RC Expand discussion on echo-state property, emphasizing that high spectral/effective dimensionality alone is insufficient for good reservoir performance. Clarify that PCA-derived $D_\text{eff}$ is not equivalent to the number of spectral comb teeth.

	\AR The discussion on the echo-state property has been expanded and we have explain why high effective dimensionality alone is not sufficient for good reservoir performance.
	Furthermore, we have also added that the PCA-derived $D_\text{eff}$ is not equivalent to the number of spectral comb teeth in page 3 of the main text.

	\RC Explain why DFT on complex amplitudes (instead of power) yields useful complementary spectral features

	\AR This information has been added in page 2 of the main text.

	\RC Briefly comment on robustness of the $\lambda_{\max}$ chaos threshold. Mark open-loop one-step-ahead prediction explicitly in Figure 1 caption to prevent misreading
	of NMSE.

	\AR We tested various values of $\lambda_{\max}$ and proved that the selected value of $\lambda_{\max}$ is robust. This has been briefly explained in page 3 of the main text and more thoroughly in Sec.~V.A.1 in the Supplemental Material.

	\RC Briefly discuss hardware practicalities including bistable branch-selection challenges in the outlook.

	\AR This has been added to the last paragraph in the main text.

	\RC Minor consistency checks for symbol formatting and leftover typos across main text and SI.

	\AR We have proofread both the main text as well as the supplemental material and fixed all typos we are aware of.

	\noindent\rule{\textwidth}{0.4pt}

	We thank both reviewers for their comments.

	Sincerely,\\
	Mahadev Sunil Kumar and Adarsh Ganesan
\end{document}
